\section{Introduction}\label{sec:introduction}

Frequency estimation affects relaying, under-frequency load shedding,
disturbance monitoring, and converter supervision in low-inertia grids. In
inverter-based-resource (IBR) systems, the estimation problem is not limited
to slow electromechanical drift. Converter-dominated dynamics compress the
disturbance timescale and can combine fast RoCoF segments, abrupt phase jumps,
harmonics, out-of-band components, and short ringdown intervals within the
same protection decision window~\cite{Milano2018LowInertia,Pinheiro2025,Ponce2025}.
Recent IBR studies also show that converter interactions can reshape the
observable frequency response itself~\cite{Bernal2025,Zoghby2024}.

These effects stress estimators through different failure modes. Phase
discontinuities can be interpreted by loop-based methods as large frequency
excursions or cycle slips. Sustained ramps stress the tracking bandwidth of
multi-cycle spectral methods. Harmonic and interharmonic content increases
spectral leakage and can bias estimators whose signal model assumes a single
dominant sinusoid. Ringdown adds a transient-recovery problem in which average
error, peak error, and trip-risk duration need not favor the same method.

The estimator families used in practice reflect this latency--robustness
tension. Windowed spectral methods such as IpDFT and Taylor--Fourier variants
reject steady harmonic content but pay observation delay~\cite{Belega2014,PlatasGarza2011}.
PLL and SOGI-based trackers are light enough for embedded deployment but can
lose phase alignment under abrupt discontinuities~\cite{Golestan2017,Wright2019}.
Kalman-filter variants can track fast transients when the state model remains
valid~\cite{Ferrero2020EKF,Chamorro2021}. Reduced-order and data-driven methods
add further cost--accuracy trade-offs, but their behavior under
relay-oriented composite stress remains less standardized~\cite{Netto2018KoopmanKF,Ting2022}.

This paper frames the contribution as a benchmark study rather than as a
single-estimator paper. The benchmark extends standard dynamic-test families
with IBR-specific harmonic, out-of-band, ringdown, and multi-event stress in a
local, non-time-synchronized setting. The study compares 16 estimators over 34
scenarios using RMSE, peak error, trip-risk duration, settling behavior, and
relative CPU cost. The contribution is therefore the disturbance-family
comparison protocol and the resulting map of trade-offs; no new estimator
family is proposed.
